\documentclass[11pt]{article}
\usepackage{fontspec}
\setmainfont{Latin Modern Roman}
\usepackage{amsmath,amssymb,mathtools,unicode-math}
\setmathfont{Latin Modern Math}
\usepackage{xurl,hyperref}
\hypersetup{colorlinks=true,urlcolor=blue,pdftitle={Problem statement}}
\newcommand{\sref}[2]{\hyperlink{source:#1:#2}{[S#2]}}
\newcommand{\cataloguescope}[1]{\paragraph{Mathematical question.}#1}
\begin{document}
\section{Algebraic Stories — Degree of the Waring map}\label{statement}

\subsection{Definitions and mathematical statement}

\paragraph{Shared notation.}$\mathbb N=\{1,2,\ldots\}$, and $\mathbb Z,\mathbb Q,\mathbb R,\mathbb C$ denote the integers, rationals, reals and complex numbers. Any other conventions are specified locally.


Let $S=\mathbb C[x,y]$ and $k\geq2$, $d\geq1$. Call $(k,d)$ perfect when $s=(kd+1)/(d+1)$ is an integer. For a perfect pair define
\[W_{k,d}:S_d^s\longrightarrow S_{kd},\qquad(g_1,\ldots,g_s)\longmapsto\sum_{j=1}^sg_j^k.
\]
Let $\widetilde W_{k,d}:S_d^s/\mathfrak S_s\to S_{kd}$ be the induced map, where the symmetric group permutes the entries. Its degree is the number of unordered tuples in the fibre over a general form. This definition uses the permutation quotient specified in the source.
\paragraph{Statement recorded in the supplied source.}
Calculate $\deg(\widetilde W_{k,d})$ for every perfect pair $(k,d)$. \sref{statement}{2}
\subsection{Short English statement}
Count the decompositions of a general binary form for every perfect pair, with the tuple counted up to permutation.
\subsection{Sources}
\begin{description}
\item[\hypertarget{source:statement:1}{S1}] ULAM.AI and the UnsolvedMath Contributors, UnsolvedMath: A Curated Collection of Open Mathematics Problems, catalogue entry. CC BY 4.0. \url{https://huggingface.co/datasets/ulamai/UnsolvedMath}.
\item[\hypertarget{source:statement:2}{S2}] R. Fröberg, S. Lundqvist, A. Oneto and B. Shapiro, \emph{Algebraic Stories from One and from the Other Pockets}, arXiv:1801.01692 (2018), §1.4, Definition 1.11 and Problem E. \url{https://arxiv.org/abs/1801.01692}.
\end{description}
\label{end:statement}
\end{document}
